\documentclass[11pt,a4paper]{article}

\usepackage[a4paper,margin=1in]{geometry}
\usepackage{hyperref}
\usepackage{listings}
\usepackage{longtable}
\usepackage{booktabs}
\usepackage{titlesec}
\usepackage{graphicx}
\usepackage{enumitem}
\usepackage{color}

\hypersetup{
	colorlinks=true,
	linkcolor=black,
	urlcolor=blue,
	pdftitle={Blentinel Technical Manual},
	pdfauthor={Blentinel Project}
}

\lstset{
	basicstyle=\ttfamily\small,
	breaklines=true,
	frame=single
}

\title{\textbf{Blentinel}\\Technical Manual}
\author{Blentinel Project}
\date{\today}

\begin{document}
	
	\maketitle
	\tableofcontents
	\newpage
	
	% ============================================================
	\section{Introduction}
	
	Blentinel is a high-performance, multi-tenant monitoring platform designed for Managed Service Providers (MSPs). It enables centralized monitoring of distributed client networks, including networks located behind NAT, private firewalls, or restricted environments.
	
	Unlike traditional monitoring systems that require inbound connectivity (VPN tunnels, open firewall ports, reverse proxies), Blentinel uses a strictly push-based model. All probe communication is outbound-initiated. This significantly reduces attack surface and simplifies deployment in remote or security-sensitive environments.
	
	Blentinel consists of two core components:
	
	\begin{itemize}
		\item \textbf{Blentinel Hub} — Central server responsible for authentication, data storage, alerting, and visualization.
		\item \textbf{Blentinel Probe} — Lightweight agent deployed inside client networks that monitors local resources and securely transmits health data to the hub.
	\end{itemize}
	
	% ============================================================
	\section{System Architecture}
	
	\subsection{High-Level Design}
	
	Blentinel uses a hub-and-spoke architecture:
	
	\begin{itemize}
		\item Probes monitor local resources.
		\item Probes initiate outbound connections to the Hub.
		\item The Hub verifies identity and stores time-series data in SQLite.
		\item Optional TLS provides transport security.
		\item Application-layer encryption provides end-to-end cryptographic guarantees.
	\end{itemize}
	
	\subsection{Data Flow}
	
	\begin{enumerate}
		\item Probe collects resource metrics.
		\item Probe signs payload using Ed25519 private key.
		\item Probe encrypts payload using X25519-derived shared secret and ChaCha20-Poly1305.
		\item Probe sends encrypted report to Hub.
		\item Hub verifies signature.
		\item Hub decrypts payload.
		\item Hub stores report in database.
	\end{enumerate}
	
	% ============================================================
	\section{Cryptographic Design}
	
	Blentinel uses layered cryptographic security:
	
	\subsection{Probe Identity (Ed25519)}
	
	Each probe generates a persistent Ed25519 keypair on first run:
	
	\begin{itemize}
		\item Private key stored in \texttt{identity.key}
		\item Public key must be registered in the Hub configuration
	\end{itemize}
	
	The Hub rejects any probe whose public key is not explicitly whitelisted.
	
	\subsection{Hub Identity (X25519)}
	
	The Hub generates an X25519 identity key:
	
	\begin{itemize}
		\item Stored in \texttt{hub\_identity.key}
		\item Public key may be configured in probe to skip handshake
	\end{itemize}
	
	\subsection{Key Exchange}
	
	If \texttt{hub\_public\_key} is not preconfigured in the probe:
	
	\begin{itemize}
		\item Probe performs initial handshake
		\item Hub provides its public key
		\item Shared secret derived via X25519
	\end{itemize}
	
	\subsection{Payload Encryption}
	
	All monitoring payloads are encrypted using:
	
	\begin{itemize}
		\item ChaCha20-Poly1305 (authenticated encryption)
		\item Ephemeral session keys derived from X25519
	\end{itemize}
	
	\subsection{Optional TLS}
	
	Transport-level HTTPS may be enabled. When enabled:
	
	\begin{itemize}
		\item Hub generates self-signed certificate
		\item Probes embed certificate for pinning
		\item Prevents MITM even with compromised CAs
	\end{itemize}
	
	% ============================================================
	\section{Installation}
	
	\subsection{Prerequisites}
	
	\begin{itemize}
		\item Rust toolchain
		\item cargo-leptos
		\item SQLite (bundled via SQLx)
	\end{itemize}
	
	\subsection{Building with blentinelmake}
	
	\begin{lstlisting}
		cargo build -p blentinelmake --release
		./target/release/blentinelmake
	\end{lstlisting}
	
	Interactive mode allows selection of hub or probe build and publish operations.
	
	\subsection{Manual Build}
	
	Hub:
	
	\begin{lstlisting}
		cargo leptos build --release
	\end{lstlisting}
	
	Probe:
	
	\begin{lstlisting}
		cargo build -p probe --release
	\end{lstlisting}
	
	% ============================================================
	\section{First-Time Setup}
	
	\subsection{Starting the Hub}
	
	On first run:
	
	\begin{itemize}
		\item \texttt{hub\_identity.key} is generated
		\item \texttt{hub\_auth.token} is generated
		\item Public hub key is displayed
	\end{itemize}
	
	Administrator must store these securely.
	
	\subsection{Starting the Probe}
	
	On first run:
	
	\begin{itemize}
		\item \texttt{identity.key} is generated
		\item Public key is displayed
	\end{itemize}
	
	Add the probe public key to Hub configuration:
	
	\begin{lstlisting}
		[[probes]]
		name = "CLIENT-SERVER"
		public_key = "HEX_PUBLIC_KEY"
	\end{lstlisting}
	
	Restart Hub and Probe.
	
	% ============================================================
	\section{Configuration}
	
	\subsection{Hub Configuration}
	
	Configuration file: \texttt{blentinel\_hub.toml}
	
	Sections include:
	
	\begin{itemize}
		\item \texttt{[server]}
		\item \texttt{[server.tls]}
		\item \texttt{[retention]}
		\item \texttt{[alerts]}
		\item \texttt{[[probes]]}
	\end{itemize}
	
	\subsection{Probe Configuration}
	
	Configuration file: \texttt{blentinel\_probe.toml}
	
	\begin{lstlisting}
		[agent]
		company_id = "HOMELAB"
		hub_url = "http://127.0.0.1:3000"
		interval = 60
		hub_public_key = "OPTIONAL"
		
		[[resources]]
		name = "Router"
		type = "ping"
		target = "192.168.1.1"
	\end{lstlisting}
	
	% ============================================================
	\section{Operational Behavior}
	
	\subsection{Hot Reload}
	
	Both Hub and Probe monitor their configuration files:
	
	\begin{itemize}
		\item Valid configuration changes applied live
		\item Invalid changes rejected
		\item Restart required only for bind-related changes
	\end{itemize}
	
	\subsection{Blackbox Database}
	
	The probe maintains \texttt{blentinel\_blackbox.db} for queued reports.
	
	If probe identity or hub public key changes:
	
	\begin{itemize}
		\item Blackbox DB is automatically reset
		\item Prevents decryption failures
	\end{itemize}
	
	% ============================================================
	\section{Alerting}
	
	Hub supports SMTP-based email alerts.
	
	Required fields when alerts enabled:
	
	\begin{itemize}
		\item SMTP server
		\item Sender address
		\item Default recipients
	\end{itemize}
	
	Thresholds configurable per metric.
	
	% ============================================================
	\section{Security Considerations}
	
	\begin{itemize}
		\item Deleting hub identity invalidates existing encrypted sessions.
		\item Deleting probe identity requires re-whitelisting.
		\item Changing hub certificate requires probe rebuild.
		\item Blackbox DB is cryptographically bound to identities.
	\end{itemize}
	
	% ============================================================
	\section{Troubleshooting}
	
	\subsection{Decryption Failed}
	
	Possible causes:
	
	\begin{itemize}
		\item Hub identity changed
		\item Probe identity changed
		\item Stale blackbox DB
	\end{itemize}
	
	\subsection{Invalid Signature}
	
	Probe not whitelisted or wrong public key.
	
	\subsection{Probe Expired}
	
	Probe not reporting within configured timeout.
	
	% ============================================================
	\section{Upgrade Procedure}
	
	\begin{enumerate}
		\item Stop services
		\item Replace binaries
		\item Preserve identity files
		\item Restart services
	\end{enumerate}
	
	% ============================================================
	\section{Directory Layout}
	
	Generated files:
	
	\begin{itemize}
		\item \texttt{hub\_identity.key}
		\item \texttt{hub\_auth.token}
		\item \texttt{hub\_tls\_cert.pem}
		\item \texttt{identity.key}
		\item \texttt{blentinel\_blackbox.db}
	\end{itemize}
	
\end{document}
